% !TEX program = xelatex
\documentclass{英语试卷}
\begin{document}
\examtitle{2024 学年第二学期浙江省精诚联盟适应性联考}{英语}
\par 高三英语 试题\par
\begin{notice}
\par 1. 答卷前，考生务必将自己的姓名、准考证号填写在答题卡上，\par
\par 2. 回答选择题时，选出每小题答案后，用铅笔把答题卡上对应题目的答案标号涂黑。如需改动，用橡皮擦干净后，再选涂其他答案标号。回答非选择题时，将答案写在答题卡上，写在本试卷上无效。\par
\par 3. 考试结束后，将本试卷和答题卡一并交回。\par
\end{notice}
\parttitle{第一部分 听力（共两节，满分30分）}
\par 做题时，先将答案标在试卷上，录音内容结束后，你将有两分钟的时间将试卷上的答案转涂到答题卡上。\par
\sectitle{第一节（共5小题：每小题1.5分，满分7.5分）}
\par 听下面 5 段对话。每段对话后有一个小题，从题中所给的 A、B、C 三个选项中选出最佳选项。听完每段对话后，你都有 10 秒钟的时间来回答有关小题和阅读下一小题。每段对话仅读一遍，\par
\question{1. What will the woman probably do in the summer vacation?}
\begin{choices}(3)
  \item Make new friends.
  \item Work in a restaurant
  \item Go to a summer camp.
\end{choices}
\question{2. What does the woman say about the bus?}
\begin{choices}(3)
  \item It arrives every 10 minutes.
  \item It passes by Ford's Theatre.
  \item It stops directly at the library.
\end{choices}
\question{3. What would the woman possibly buy?}
\begin{choices}(3)
  \item Skimmed milk
  \item Organic products.
  \item Whole grain bread.
\end{choices}
\question{4. What is the man doing?}
\begin{choices}(3)
  \item Exploring a museum. 微信公众号 IAI English
  \item Visiting an exhibition.
  \item Making a consultation.
\end{choices}
\question{5. What are the speakers going to do together?}
\begin{choices}(3)
  \item Have lunch downtown.
  \item Meet at the dentist's.
  \item Attend a work meeting.
\end{choices}
\sectitle{第二节（共 15 小题：每小题 1.5 分，满分 22.5 分）}
\par 听下面 5 段对话或独白。每段对话或独白后有几个小题，从题中所给的 A、B、C 三个选项中选出最佳选项。听每段对话或独白前，你将有时间阅读各个小题，每小题 5 秒钟；听完后，各小题将给出 5 秒钟的作答时间。每段对话或独白读两遍。\par
\par 听第 6 段材料，回答第 6、7 题，\par
\question{6. What do we know about Jasmine?}
\begin{choices}(3)
  \item She's in need of attention and care.
  \item She keeps a dog for companionship.
  \item She is considering having a dog as a pet.
\end{choices}
\question{7. What is Alex's advice to Jasmine?}
\begin{choices}(3)
  \item Receive training.
  \item Be committed to the dog.
  \item Do research on the breeds.
\end{choices}
\par 听第 7 段材料，回答第 8 至 10 题，\par
\question{8. Why does Alan want to help around the house?}
\begin{choices}(3)
  \item To live up to parents' expectations
  \item To make money for Dad's birthday gift.
  \item To share the housework equally with Mom.
\end{choices}
\question{9. What's Alan likely to do next?}
\begin{choices}(3)
  \item Do the laundry.
  \item Look for more chores.
  \item Buy a watch with Dad.
\end{choices}
\question{10. How does the woman feel now?}
\begin{choices}(3)
  \item Valued.
  \item Delighted.
  \item Concerned.
\end{choices}
\par 听第 8 段材料，回答第 11 至 13 题，\par
\question{11. Where does the conversation probably take place?}
\begin{choices}(3)
  \item At a square.
  \item At a tube station.
  \item At a ticket office.
\end{choices}
\question{12. How frequently do the trains run during rush hours?}
\begin{choices}(3)
  \item Every 15 minutes.
  \item Every 20 minutes.
  \item Every 5 to 10 minutes.
\end{choices}
\question{13. What will the man probably do next?}
\begin{choices}(3)
  \item Buy a weekly pass.
  \item Check his departure time
  \item Take the yellow line train.
\end{choices}
\par 听第 9 段材料，回答第 14 至 17 题。\par
\question{14. What does the woman like?}
\begin{choices}(3)
  \item Handcrafted items.
  \item Various handbags.
  \item Trendy accessories.
\end{choices}
\question{15. How did the woman get her handbag?}
\begin{choices}(3)
  \item She made it by herself.
  \item She bought it during her trip.
  \item It was a gift from local people.
\end{choices}
\question{16. What is required to make a bag?}
\begin{choices}(3)
  \item A large workforce.
  \item Time and patience.
  \item Fibers of various plants.
\end{choices}
\question{17. What makes the handbag special?}
\begin{choices}(3)
  \item Its unusual size.
  \item Its unique appearance.
  \item Its traditional material.
\end{choices}
\par 听第 10 段材料，回答第 18 至 20 题，\par
\question{19. How did the dolphins help the speaker?}
\begin{choices}(3)
  \item By pulling him underwater.
  \item By keeping him near the shore.
  \item By playing and supporting him.
\end{choices}
\question{20. What is the speaker's main message?}
\begin{choices}(3)
  \item Facing fears leads to growth.
  \item Dolphins can cure human fears.
  \item Safety promotes personal growth.
\end{choices}
\parttitle{第二部分 阅读理解（共两节，满分 50 分）}
\sectitle{第一节（共15小题：每小题2.5分，满分37.5分）}
\par 阅读下列短文，从每题所给的 A、B、C 和 D 四个选项中，选出最佳答案，\par
\passagelabel{A}
\par College of Arts and Sciences at UP\par
\par The oldest and largest school at University of Portland (UP), the College of Arts and Sciences offers degrees in the humanities, natural sciences, performing arts, social sciences, and interdisciplinary fields. A liberal arts and sciences degree provides a versatile (灵活多向的) pathway to a wide range of careers, and the College takes full advantage of UP's ideal location in creating hands-on career-building experiences out in the field, whichever one you choose!\par
\par Popular Programs\par
\par Environmental Science/ Environmental Ethics \& Policy\par
\par Psychology\par
\par Pre-Health\par
\par Political Science and Pre-Law\par
\par Gender, Women, and Sexuality Studies\par
\par Ethnic Studies\par
\par Humanities\par
\par The Portland Advantage\par
\par With close proximity (临近) to the outdoors, government agencies, leading medical institutions, arts and culture venues, and more, UP gives students a great variety of opportunities to connect with the community and gain valuable hands-on experiences in a wide range of careers.\par
\par R1-Level Research\par
\par UP goes head-to-head with Research 1 (R1) institutions, with unmatched opportunities for hands-on research work with faculty (全体教员) members throughout the academic year and summer in disciplines as varied as theology, biology, ethnic studies, and mathematics.\par
\par Hands-on Connections\par
\par While R1 universities increasingly rely on graduate students as instructors, at UP all classes are taught by our professional and admirable faculty. That means your professors will really get to know you and will provide expert guidance, advice, and unique academic opportunities!\par
\question{21. What can we learn about College of Arts and Sciences?}
\begin{choices}(2)
  \item It is located in the countryside.
  \item It is the most popular school at UP.
  \item It seeks employment for its students.
  \item It prepares its students for various careers.
\end{choices}
\question{22. What advantage does UP have over R1 universities?}
\begin{choices}(2)
  \item Its disciplines are more varied.
  \item Its professors are more admirable.
  \item It offers more tailored learning experiences.
  \item It provides more chances for hands-on research.
\end{choices}
\question{23. Where is this text probably taken from?}
\begin{choices}(2)
  \item A course plan.
  \item A program proposal.
  \item An admission brochure.
  \item An educational report.
\end{choices}
\passagelabel{B}
\par A few years ago, a chance encounter with traditional Chinese medicine (TCM) excited my deep interest in this ancient culture. Greatly fascinated, I decided to pursue a systematic study.\par
\par I applied to Yunnan University of Chinese Medicine, and I was overjoyed the day I received my admission letter. At the beginning, the ancient Chinese classics and complex processes of herb collection, identification and combination were quite challenging to me. Fortunately, I had the patient guidance of professors and the enthusiastic help of classmates to rely on.\par
\par I remember when I first started learning acupuncture (针灸), I was often too hesitant to actually use the needles. To encourage me, my senior classmates generously acted as patients, braving the discomfort to let me practice acupuncture techniques on their bodies. With their support and motivation, my acupuncture skills steadily improved, and my studies quickly got on track.\par
\par During my learning journey, the most memorable experience was my internship (实习) at the Yunnan Provincial Hospital of Traditional Chinese Medicine. There, I had the opportunity to participate directly in patient care. I looked forward to every Saturday afternoon when I could shadow my teacher. After each patient consultation, she would give me a detailed explanation: what the patient's pulse (脉搏) indicated, what it represented in terms of their conditions, and the role of each herb in the prescription (处方). This approach helped shape my diagnostic thinking and was incredibly beneficial to me. Even though this meant she could never leave work on time, she never complained. As an ancient Chinese saying goes, "It’s easier to get a teacher of classics than a teacher of character."\par
\par My every day in China has been filled with challenges and surprises. TCM has taught me not just the skills to heal and save lives but how to view life. I have come to understand that the power of culture lies more in its ability to connect hearts than in its uniqueness. People from different cultural backgrounds can actually form strong connections.\par
\question{24. Which of the following did the author find difficult when starting studying TCM?}
\begin{choices}(2)
  \item How to collect ancient Chinese classics.
  \item How to combine different herbs properly.
  \item How to communicate effectively with professors.
  \item How to understand the ancient Chinese philosophy
\end{choices}
\question{25. Why does the author mention the experience of learning acupuncture?}
\begin{choices}(2)
  \item To show the complexity of acupuncture.
  \item To introduce the procedures of studying TCM.
  \item To illustrate the supportive learning environment.
  \item To stress the dangers of improper acupuncture techniques.
\end{choices}
\question{26. What does the author mean by the saying in paragraph 4?}
\begin{choices}(2)
  \item His TCM teacher is highly professional.
  \item His TCM teacher is selfless and devoted.
  \item His TCM teacher is strict and demanding.
  \item His TCM teacher is easy to get along with.
\end{choices}
\question{27. What does the author learn from the experience of studying TCM?}
\begin{choices}(2)
  \item Life is precious.
  \item Hard work pays off.
  \item Uniqueness unites, not divides.
  \item Culture is a bridge, not a wall.
\end{choices}
\passagelabel{C}
\par You've seen it happen: you present a fact — clear, undeniable — and instead of reconsidering their position, the other person doubles down. "I just don't believe it," they say. This confusing resistance isn't just frustrating; it's deeply human. Our brains are wired not merely to seek truth, but to defend identity. Welcome to the mind's firewall—our internal system that protects beliefs from intruding (侵入的) evidence. And yet, understanding this firewall is the first step to disarming it, for ourselves and others.\par
\par In Elizabeth Kolbert's article "Why Facts Don’t Change Our Minds," she draws on studies in cognitive science to show that reasoning evolved not to find truth, but to win arguments. Our brains often favor "motivated reasoning"— seeking information that supports our views and discounting evidence that challenges them. This tendency, known as "confirmation bias (偏见)”, helps us maintain a consistent sense of self and stay in line with our social groups. This makes conversations less about learning and more about defending one’s territory.\par
\par Think of your mind like a computer. Facts are like data packets trying to get through your firewall. But unless they meet certain criteria — emotional relevance, social safety, identity compatibility (兼容) — they're blocked. The mind doesn't simply take in raw information; it tries to make sense of it. Facts alone aren't enough to truly connect or stick. This is especially true in emotionally charged fields like politics, health, or identity. When we feel threatened, the emotional brain overrides the rational (理智的) brain. What feels true outweighs what is true.\par
\par If facts can't get past the firewall, what can? As discussed in my book The Art of Change, it's not about tearing down the firewall, but about finding alternative paths around it — narrative, emotion, and relationship.\par
\par Next time you're trying to change someone's mind — or your own — remember: lead with connection, not argument. Start with shared experiences or values. Use stories, not just analytics. And stay open yourself. Try this: Think of a belief you hold strongly. Then write down three reasons someone might see it differently — not arguments, but genuinely thoughtful reasons. Doing this won't make you weaker in your beliefs. It'll make you wiser in your understanding.\par
\question{28. What does the author think of the mind's firewall?}
\begin{choices}(2)
  \item It is natural.
  \item It is confusing.
  \item It is beneficial.
  \item It is necessary.
\end{choices}
\question{29. Why do our brains favor "motivated reasoning"?}
\begin{choices}(2)
  \item It helps people seek truth.
  \item It allows people to defend identity.
  \item It makes people learn more about themselves.
  \item It enables people to overcome any challenges.
\end{choices}
\question{30. What is paragraph 3 mainly about?}
\begin{choices}(2)
  \item The scientific method for processing facts.
  \item The underlying logic of the mind's firewall.
  \item The biological structure of the human brain.
  \item The conflict between emotional and rational brains.
\end{choices}
\question{31. What does the author suggest readers do to change others' minds?}
\begin{choices}(2)
  \item Argue against their beliefs.
  \item Present a through analysis to them.
  \item Form an emotional bond with them.
  \item Train them to be more open-minded.
\end{choices}
\passagelabel{D}
\par Amir Goldberg discovered a funny thing about American culture on his way to researching the complex topic. Goldberg and his colleagues at Stanford and Yale universities analyzed millions of Yelp and Netflix reviews to reveal that people considered the most culturally adventurous are actually the most resistant to experiences perceived as "crossing the line."\par
\par That is, those described as "cultural omnivores (杂食动物)” — because they eat Thai for lunch, play bocce ball after work, and stream a French film that night — are the very ones opposed to mixing it up. No hummus on their hot dogs, forget about spaghetti Westerns, and do not mention Switched-On Bach. Those offerings are not considered culturally authentic. They are a hodgepodge (混杂物) to which these folks would likely wrinkle their collective noses.\par
\par "We find these people hate the most atypical offerings," says Goldberg. "They can pretend to be the most open, but it turns out they are not. By being multicultural, they are the most conservative and the most resistant to changes."\par
\par The findings are opposite to previous research, which viewed omnivorousness as synonymous with openness. Instead, Goldberg, along with Michael T. Hannan, professor emeritus of organizational behavior at Stanford GSB and Balazs Kovacs of Yale University, write in their paper that "a taste for variety decreases one's receptivity to cultural innovation."\par
\par But as the team's research points out, there is a subset of consumers who are open to change. "For cultural novelty to have an impact and occasionally drive change, some agents must be tolerant to breaking traditional boundaries," he says.\par
\par The trick is to find and develop them, as Steve Jobs did with the iPhone. Goldberg notes that when the iPhone was first introduced, it confused the public. What is it? A phone? A computer? Many saw the product and dismissed it, deciding phones are for personal communication and computers are for the office. The professor notes that it was Jobs and his loyal followers who changed minds, changed habits, and ultimately changed the global culture.\par
\question{32. What do we know about "cultural omnivores"?}
\begin{choices}(2)
  \item They are unwilling to change.
  \item They are adventurous customers.
  \item They prioritize American culture.
  \item They hate culturally authentic offerings.
\end{choices}
\question{33. What does the underlined phrase "synonymous with" in paragraph 4 mean?}
\begin{choices}(2)
  \item Heavily dependent on.
  \item Greatly superior to.
  \item Closely connected with.
  \item Completely contrary to
\end{choices}
\question{34. What does Goldberg think of the iPhone?}
\begin{choices}(2)
  \item Confusing.
  \item Pioneering.
  \item Convenient.
  \item Effective.
\end{choices}
\question{35. What is the key to overcoming cultural resistance according to the text?}
\begin{choices}(2)
  \item Educating the public.
  \item Setting traditional boundaries.
  \item Developing innovative products.
  \item Making use of the open-minded.
\end{choices}
\sectitle{第二节（共10小题；每小题1.5分，满分15分）}
\par 阅读下面短文，从短文后的选项中选出能填入空白处的最佳选项。选项中有两项为多余选项。\par
\par Dealing with negativity can be difficult, especially when you're not happy with who you are. \blanknum{36} However, laughing at yourself is a surprisingly effective way to deal with negative circumstances. Here are some tips to learn to laugh at yourself more in life.\par
\par Accept yourself. You know that nobody is perfect. \blanknum{37} When talking about the "negatives", it's extremely important to know which parts of yourself can be changed and which can't. For example, there's no point in being angry at yourself because you aren't as tall as you'd like to be. You can try all you want, but you'll never change those features of yourself.\par
\par \blanknum{38} Accepting yourself is often so hard because you aren't even looking at yourself — you're just looking at how well you do a certain thing. You've probably done it yourself. You didn't lift as much weight in the gym as the guy next to you? You're weak. You didn't finish all of the work you planned to do? You're lazy. \blanknum{39} So instead, you should realize that a single action doesn't determine who you are. This means that laughing about a mistake you once made doesn't mean you're laughing at who you are as a person.\par
\par Be humble. You can't laugh at yourself while at the same time trying to maintain the image that you're awesome at everything you do. In other words, you'd better avoid showing off and be more humble instead. A good tip is to think of how insignificant you are in the universe. \blanknum{40} It reminds you that whatever problems you're dealing with right now don't matter in the end.\par
\begin{choices}(1)
  \item Be mindful of others.
  \item This knowledge can be freeing in a way
  \item Separate yourself from your performance.
  \item So how can you learn to laugh about yourself?
  \item So why would you expect yourself to be any different?
  \item You can make a mistake, but that doesn't make you the mistake.
  \item The last thing you want at this point is for someone to laugh at you.
\end{choices}
\parttitle{第三部分 语言知识运用（共两节，满分 30 分）}
\sectitle{第一节（共15小题；每小题1分，满分15分）}
\par 阅读下面短文，从短文后各题所给的 A、B、C、D 四个选项中，选出可以填入空白处的最佳选项，\par
\par My grandmother set up a very successful business as a widowed (寡居的) woman in the South in the 1950s. She had to be \blanknum{41} and my father was too — handsome and ambitious, to be sure, but also tough. He prided himself on the \blanknum{42} of his children. Even as a little girl \blanknum{43} to jump rope, I knew I could \blanknum{44} approval through physical trials.\par
\par His \blanknum{45} was as intense as his expectations were high. The absence in his life that his \blanknum{46} had left behind in his life was vast. So he was determined to be present, to never \blanknum{47} so much as a Little League game. He was home every day at five thirty, \blanknum{48} as clockwork. He coached my soccer team in the park. He had no \blanknum{49} of being a father, because fatherhood had never been modeled for him. \blanknum{50}, he figured it out somehow.\par
\par From my father, I learned what it \blanknum{51} to be a businessperson. I understood that his employees needed to see his \blanknum{52} On weekends and holidays, he would drive around and \blanknum{53} the clerks to make sure business was running smoothly. And on snow days a few times a year, I would go with my dad to help \blanknum{54} the stores so that people could get food, essentials, and gas. We took pride in \blanknum{55} a public service; we were part of the fabric of our town.\par
\cz{41}{warm-hearted}{strong-willed}{open-minded}{good-tempered}
\cz{42}{kindness}{independence}{achievements}{talents}
\cz{43}{learning.}{promising}{failing}{pretending}
\cz{44}{win}{express}{require}{deny}
\cz{45}{passion}{ambition}{curiosity}{devotion}
\cz{46}{mother}{father}{grandmother}{child}
\cz{47}{interrupt}{escape}{miss}{forget}
\cz{48}{predictable}{accessible}{replaceable}{respectable}
\cz{49}{intention}{concept}{chance}{fear}
\cz{50}{Instead}{Otherwise}{Therefore}{Nevertheless}
\cz{51}{took}{offered}{meant}{changed}
\cz{52}{generosity}{courage}{leadership}{confidence}
\cz{53}{fill in for}{check in on}{take care of}{cut in on}
\cz{54}{promote}{empty}{clean}{open}
\cz{55}{providing}{expanding}{using}{restoring}
\sectitle{第二节（共10小题；每小题1.5分，满分15分）}
\par 阅读下面短文，在空白处填入 1 个适当的单词或括号内单词的正确形式.\par
\par Chinese peonies (牡丹) are not just beautiful flowers; they also hold deep symbolic meaning in Chinese culture.\par
\par Often \blanknum{56} (refer) to as the King of Flowers, the peony is one of the most culturally significant flowers in China, with a history dating back over 2,000 years. It is believed \blanknum{57} (represent) wealth, prosperity (繁荣), and honor.\par
\par One of the reasons for the peony's cultural significance is its historical \blanknum{58} (associate) with the imperial family. During the Tang Dynasty, the peony \blanknum{59} (consider) the official flower of the imperial court and only the imperial family and high-ranking \blanknum{60} (official) could plant and enjoy these delicate flowers.\par
\par In addiction, the peony was deeply rooted in Chinese mythology. According to Chinese legend, the peony is associated with a beautiful goddess named Xi Wangmu, \blanknum{61} planted peonies in her garden to symbolize her immortality (永生). Thus, the peony is often associated with longevity and the \blanknum{62} (pursue) of immortality.\par
\par Also, the peony holds \blanknum{63} special place in Chinese weddings. It is greatly admired \blanknum{64} a symbol of love and happiness and is often used in wedding ceremonies and decorations.\par
\par In conclusion, the cultural significance of Chinese peonies cannot be underestimated. Over centuries, their rich history, deep symbolism, and aesthetic beauty \blanknum{65} (make) them one of the most treasured flowers in China.\par
\parttitle{第四部分 写作（共两节，满分 40 分）}
\sectitle{第一节（满分15分）}
\par 假定你是李华。上周你校举办了一场以“环境科学”为主题的科学博览会。你的外国笔友 Andy 对此很感兴趣，写信向你了解情况。请你给他回信，内容包括：\par
\par (1) 活动目的和内容;\par
\par (2) 你的感受。\par
\par 注意:\par
\par Dear Andy, Best regards, Li Tua\par
\par (1) 写作词数应为 80 左右;\par
\par (2) 可适当增加细节，以使行文连贯，\par
\sectitle{第二节（满分25分）}
\par 阅读下面材料，根据其内容和所给段落开头语续写两段，使之构成一篇完整的短文。\par
\par Here are a few things you need to know about my dad. First, he was a slightly impatient kind of a guy. If he wanted something done, he wanted it done NOW, if not sooner. Second, you need to know that Dad was an early riser. Every single day. No exceptions. No excuses. Work or no work. Holidays or vacation days, there was my father, up and ready to go. Oh, and you should know that Mom was a night owl, up late watching a late movie or a wrestling natch.\par
\par One more thing about my dad. He was color-blind. That usually doesn't matter. And he was a fairly conservative kind of person. For instance, he bought our new home out in the country and objected immediately to the stucco (灰尼) being a pale peach color. Way too "flashy" for Dad. He hired a painter to cover up all that peach with a nice, not-flashy gray. Dad was fond of gray anyway.\par
\par And so, we arrived at the Saturday morning when Dad sprang out of bed as bright-eyed as ever, felt restless, and looked around for something that needed doing. Mom was sound asleep, as usual. My brother and I slept in a bit. There was nobody around to keep an eye on Dad and whatever he decided to do next.\par
\par So, Dad decided that the porch (门廊) on our newly gray house needed a fresh coat of paint. He headed off to the hardware store to get paint and painting supplies. He did not tell any of us. Nor did he invite any of us along.\par
\par We've never been sure what happened at the hardware store. Did he ask for advice? Did he study color samples? Did another shopper help him choose? Who knows? But we're all pretty sure that Dad had no clue what color he was painting the wooden porch. My brother and I woke up, wondered where Dad was, and went looking for him. 注意:\par
\par (1) 续写词数应为 150 左右;\par
\par (2) 请按如下格式在答题卡的相应位置作答，\par
\writingpara{Paragraph 1: We froze in our tracks when we stepped outside. \blanknum{}}
\writinglines{8}
\writingpara{Paragraph 2: Looking at the bright pink porch, we three exchanged looks, unsure how to break the news. \blanknum{}}
\writinglines{8}
\end{document}
